\documentclass[10pt,a4paper]{amsart}
\usepackage[margin=19mm]{geometry}
\usepackage{amsmath,amssymb,mathtools}
\usepackage[scr=rsfs,cal=euler]{mathalfa}
\usepackage{xfrac}
\usepackage[unicode,hidelinks,bookmarks=false]{hyperref}
\newcommand{\ie}{i.e.\ }
\newcommand{\cf}{cf.\ }
\newcommand{\resp}{respectively\ }
\newcommand{\Subsection}{\subsection}
\providecommand{\nobreakdash}{\nobreak\hbox{-}}
\setlength{\emergencystretch}{2em}
\multlinegap0pt
\usepackage{bm}
\usepackage{bbm}
\usepackage{dsfont}
\usepackage{xspace}
\usepackage{cancel}
\usepackage{mathtools}
\usepackage{amsthm}
\makeatletter
\@ifundefined{theorem}{\newtheorem{theorem}{Theorem}}{}
\makeatother
\title[Ricci curvature bounds for Statistical Submersions]{Ricci curvature bounds for Statistical Submersions}
\author{Ravindra Singh}
\date{Source: arXiv:2608.21760v1 (CC BY 4.0)}
\begin{document}
\maketitle

\noindent\emph{Source and licence.} Ricci curvature bounds for Statistical Submersions. Version arXiv:2608.21760v1 (CC BY 4.0), distributed under the Creative Commons Attribution 4.0 International licence, as recorded in the arXiv licence metadata for this exact version and shown on the abstract page https://arxiv.org/abs/2608.21760v1. Full rights notice: licenses/arxiv-2608.21760-CC-BY-4.0.txt.\par\smallskip

\noindent\textit{Editorial scope.} Counted unit: the Theorem whose source label is \texttt{eq-CR-SS} in the article (source0.tex lines 3556--3599, source label \texttt{eq-CR-SS}), delivered together with the article's own complete original proof of it (source0.tex lines 3601--3696, one \texttt{proof} environment of 96 lines). No mathematics is written, repaired or re-derived: statement and proof are the article's own text line for line. Required context retained uncounted: eq-nrm-A (lines 1312--1322); eq-Gauss-hor (lines 1452--1468); eq-Ric-hor (lines 1528--1537). Every definition, display and label of those lines is retained unchanged. Omitted from this package and not redistributed: 1--1311, 1323--1451, 1469--1527, 1538--3555, 3600--3600, 3697--4255. Nothing inside a retained excerpt is omitted or altered: every retained body is delivered exactly as the article's lines are written, so no omission receipt of any kind accompanies this package. Citations: the retained excerpts carry no citation command at all, so no bibliography entry of the article is transcribed and no reference is redistributed. Reference adaptation: none was needed and none was made; every reference inside the delivered unit resolves inside it, so no unresolved reference or citation is produced. Printed slips kept literal: no printed word, symbol, hypothesis or number of the counted statement or of its proof was altered. Layout and class: the article's own class is \texttt{article} with packages including geometry, amsmath, amssymb, amsfonts, amsthm, mathrsfs, graphicx, booktabs, multirow, xcolor, enumerate, appendix, hyperref; the wrapper is the lane's pinned \texttt{amsart} editorial wrapper at 10pt on A4, which restates the theorem-like environments the retained text needs and adds the article's own macro definitions that the retained text needs (none), with extra wrapper packages (bm, bbm, dsfont, xspace, cancel, mathtools). The delivered numbering is the unit's own. Assets: no retained span contains a figure environment, a graphics-inclusion command, a PDF-page inclusion, a table or a TikZ picture; no image file is copied into this package, the package declares no asset dependency and no image file is redistributed with it. Licence: version arXiv:2608.21760v1 of this article is distributed under the Creative Commons Attribution 4.0 International licence, as recorded in the arXiv licence metadata for this exact version and shown on the abstract page https://arxiv.org/abs/2608.21760v1. A full rights notice is delivered at licenses/arxiv-2608.21760-CC-BY-4.0.txt. Characters: every retained excerpt is pure ASCII, so the delivered engine, XeLaTeX with its default Latin Modern text font, reports no missing character.
\begin{equation}
\left\|\operatorname{trace}A\right\|^{2}
=
\|\sigma\|^{2}
=
\sum_{\alpha=1}^{r}
\left(
\sum_{j=1}^{s}A_{jj}^{\alpha}
\right)^{2},
\label{eq-nrm-A}
\end{equation}

\begin{align}
R^{M}(U_1,U_2,U_3,U_4)
={}&
R^{\mathcal H}(U_1,U_2,U_3,U_4)
-
g_1(A_{U_2}U_3,A^{*}_{U_1}U_4)
\nonumber\\
&+
g_1\left(
(A_{U_1}+A^{*}_{U_1})U_2,
A^{*}_{U_3}U_4
\right)
\nonumber\\
&+
g_1(A_{U_1}U_3,A^{*}_{U_2}U_4).
\label{eq-Gauss-hor}
\end{align}

\begin{equation}
\operatorname{Ric}_{\mathcal H}^{M}(h_1)
=
\sum_{j=2}^{s}
R^{M}(h_1,h_j,h_j,h_1), \quad \operatorname{Ric}^{\mathcal H}(h_1)
=
\sum_{j=2}^{s}
R^{\mathcal H}(h_1,h_j,h_j,h_1).
\label{eq-Ric-hor}
\end{equation}

\begin{theorem}
Let
\[
F:(M,\nabla ,g_{1})\longrightarrow (N,\nabla ^2,g_{2})
\]
be a statistical submersion between two statistical manifolds. Then, for
every unit horizontal vector $h_{1}\in\mathcal H_{p}$,
\begin{align}
\operatorname{Ric}_{\mathcal H}^{M}(h_{1})
\leq{}&
\operatorname{Ric}^{\mathcal H}(h_{1})
+\frac14\|\operatorname{trace}\mathcal A\|^{2}
+\frac18\sum_{\alpha=1}^{r}\sum_{j=2}^{s}
(\mathcal A_{1j}^{*\alpha})^{2}
\nonumber\\
&-2\sum_{\alpha=1}^{r}\sum_{j=2}^{s}
\left(
\mathcal A_{1j}^{\alpha}
+\frac14\mathcal A_{1j}^{*\alpha}
\right)^{2}.
\label{eq-CR-SS}
\end{align}
Consequently,
\begin{equation}
\operatorname{Ric}_{\mathcal H}^{M}(h_{1})
\leq
\operatorname{Ric}^{\mathcal H}(h_{1})
+\frac14\|\operatorname{trace}\mathcal A\|^{2}
+\frac18\sum_{\alpha=1}^{r}\sum_{j=2}^{s}
(\mathcal A_{1j}^{*\alpha})^{2}.
\label{eq-CR-SS-1}
\end{equation}
Equality in \eqref{eq-CR-SS-1} holds if and only if
\[
\mathcal A_{11}^{\alpha}
=
\sum_{j=2}^{s}\mathcal A_{jj}^{\alpha},
\qquad
\mathcal A_{1j}^{\alpha}
=
-\frac14\mathcal A_{1j}^{*\alpha},
\]
for every $\alpha=1,\ldots,r$ and $j=2,\ldots,s$.
\end{theorem}

\begin{proof}
From \eqref{eq-Gauss-hor} and the definition of the horizontal Ricci
curvature \eqref{eq-Ric-hor}, we obtain
\begin{align}
\operatorname{Ric}_{\mathcal H}^{M}(h_{1})
={}&
\operatorname{Ric}^{\mathcal H}(h_{1})
+\sum_{\alpha=1}^{r}
\mathcal A_{11}^{\alpha}
\sum_{j=2}^{s}\mathcal A_{jj}^{\alpha}
\nonumber\\
&-2\sum_{\alpha=1}^{r}\sum_{j=2}^{s}
(\mathcal A_{1j}^{\alpha})^{2}
-\sum_{\alpha=1}^{r}\sum_{j=2}^{s}
\mathcal A_{1j}^{\alpha}\mathcal A_{1j}^{*\alpha}.
\label{eq-hor-1}
\end{align}
Using the identity
\[
-2a^{2}-ab
=
\frac18b^{2}
-2\left(a+\frac14b\right)^{2},
\]
equation \eqref{eq-hor-1} becomes
\begin{align}
\operatorname{Ric}_{\mathcal H}^{M}(h_{1})
={}&
\operatorname{Ric}^{\mathcal H}(h_{1})
+\sum_{\alpha=1}^{r}
\mathcal A_{11}^{\alpha}
\sum_{j=2}^{s}\mathcal A_{jj}^{\alpha}
\nonumber\\
&+\frac18\sum_{\alpha=1}^{r}\sum_{j=2}^{s}
(\mathcal A_{1j}^{*\alpha})^{2}
\nonumber\\
&-2\sum_{\alpha=1}^{r}\sum_{j=2}^{s}
\left(
\mathcal A_{1j}^{\alpha}
+\frac14\mathcal A_{1j}^{*\alpha}
\right)^{2}.
\label{eq-hor-3}
\end{align}
Next, applying the elementary inequality
\[
d_{1}d_{2}
\leq
\frac{(d_{1}+d_{2})^{2}}{4},
\]
we obtain
\begin{equation}
\mathcal A_{11}^{\alpha}
\sum_{j=2}^{s}\mathcal A_{jj}^{\alpha}
\leq
\frac14
\left(
\mathcal A_{11}^{\alpha}
+\sum_{j=2}^{s}\mathcal A_{jj}^{\alpha}
\right)^{2}.
\label{eq-hor-4}
\end{equation}
Substituting \eqref{eq-hor-4} into \eqref{eq-hor-3} and using
\eqref{eq-nrm-A}, we arrive at
\begin{align}
\operatorname{Ric}_{\mathcal H}^{M}(h_{1})
\leq{}&
\operatorname{Ric}^{\mathcal H}(h_{1})
+\frac14\|\operatorname{trace}\mathcal A\|^{2}
+\frac18\sum_{\alpha=1}^{r}\sum_{j=2}^{s}
(\mathcal A_{1j}^{*\alpha})^{2}
\nonumber\\
&-2\sum_{\alpha=1}^{r}\sum_{j=2}^{s}
\left(
\mathcal A_{1j}^{\alpha}
+\frac14\mathcal A_{1j}^{*\alpha}
\right)^{2},
\label{eq-hor-5}
\end{align}
which proves \eqref{eq-CR-SS}. Since the last term in \eqref{eq-hor-5} is non-positive, inequality
\eqref{eq-CR-SS-1} follows immediately.

Finally, equality in \eqref{eq-CR-SS-1} holds if and only if equality is
attained in \eqref{eq-hor-4} and the squared term in
\eqref{eq-hor-5} vanishes identically. Hence
\[
\mathcal A_{11}^{\alpha}
=
\sum_{j=2}^{s}\mathcal A_{jj}^{\alpha},
\qquad
\mathcal A_{1j}^{\alpha}
=
-\frac14\mathcal A_{1j}^{*\alpha},
\]
for every $\alpha=1,\ldots,r$ and $j=2,\ldots,s$. This completes the
proof.
\end{proof}

\end{document}
